%=======================================================================
%  Catchment Analyzer -- Phase 1 Report
%  Compile:  pdflatex report.tex  (run twice for the table of contents)
%=======================================================================
\documentclass[11pt,a4paper]{article}

%----------------------------------------------------------- packages
\usepackage[utf8]{inputenc}
% OT1 deliberately, not T1: this TeX Live install ships Computer Modern as
% scalable Type 1 outlines (cmr10.pfb) but the T1-encoded EC fonts only as
% METAFONT bitmaps. OT1 therefore gives a sharper PDF and lets microtype's
% font expansion work. If you build where lmodern is installed, you may
% switch to \usepackage[T1]{fontenc}\usepackage{lmodern} for better output.
\usepackage[a4paper,margin=25mm,headheight=15pt]{geometry}
\usepackage{amsmath}
\usepackage{amssymb}
\usepackage{graphicx}
\usepackage{booktabs}
\usepackage{array}
\usepackage{longtable}
\usepackage{colortbl}
\usepackage{xcolor}
\usepackage{enumitem}
\usepackage{titlesec}
\usepackage{fancyhdr}
\usepackage{float}
\usepackage{caption}
\usepackage[protrusion=true,expansion=false]{microtype}
\usepackage{listings}
\usepackage{framed}
\usepackage[hidelinks]{hyperref}

%----------------------------------------------------------- metadata
\newcommand{\repourl}{https://github.com/devreddy1015/catchment\_analyzer}
\newcommand{\repourlraw}{https://github.com/devreddy1015/catchment_analyzer}
% >>> EDIT THIS ONE LINE with your deployed host before submitting <<<
\newcommand{\apihost}{http://10.1.75.53:5249}
\newcommand{\apiroute}{\apihost/api/analyzeContour}

%----------------------------------------------------------- palette
\definecolor{ink}{HTML}{1B2B34}
\definecolor{water}{HTML}{1F6F8B}
\definecolor{terrain}{HTML}{6B8E4E}
\definecolor{clay}{HTML}{B5651D}
\definecolor{mist}{HTML}{F4F6F7}
\definecolor{rule}{HTML}{C8D2D6}
\definecolor{kw}{HTML}{1F6F8B}
\definecolor{cmt}{HTML}{7A8B91}
\definecolor{str}{HTML}{8A5A2B}

\hypersetup{
  colorlinks=true, linkcolor=water, urlcolor=water, citecolor=water,
  pdftitle={Catchment Analyzer -- Phase 1 Report},
  pdfauthor={Aman Pratap}
}

%----------------------------------------------------------- headings
\titleformat{\section}{\Large\bfseries\color{ink}}{\thesection}{0.7em}{}
\titleformat{\subsection}{\large\bfseries\color{water}}{\thesubsection}{0.7em}{}
\titleformat{\subsubsection}{\normalsize\bfseries\color{ink}}{\thesubsubsection}{0.6em}{}
\titlespacing*{\section}{0pt}{18pt}{8pt}

\pagestyle{fancy}
\fancyhf{}
\renewcommand{\headrulewidth}{0.4pt}
\lhead{\small\color{ink}Catchment Analyzer}
\rhead{\small\color{ink}Phase 1 Report}
\cfoot{\small\thepage}

\captionsetup{font=small,labelfont={bf,color=ink}}
\setlist[itemize]{leftmargin=1.3em,itemsep=2pt,topsep=3pt}
\setlist[enumerate]{leftmargin=1.5em,itemsep=2pt,topsep=3pt}

%----------------------------------------------------------- listings
\lstdefinelanguage{pseudo}{
  morekeywords={function,procedure,return,for,foreach,while,if,else,elif,then,do,
    end,input,output,let,in,to,not,and,or,break,continue,repeat,until,new,push,
    pop,append,assert,raise,where,each,of,is,from,by,with,all,any,sort,descending},
  sensitive=true,
  morecomment=[l]{//},
  morestring=[b]",
}

\lstdefinestyle{pseudostyle}{
  language=pseudo,
  basicstyle=\ttfamily\small\color{ink},
  keywordstyle=\color{kw}\bfseries,
  commentstyle=\color{cmt}\itshape,
  stringstyle=\color{str},
  numbers=left, numberstyle=\tiny\color{cmt}, numbersep=9pt,
  showstringspaces=false,
  breaklines=true, breakatwhitespace=true,
  columns=fullflexible, keepspaces=true,
  frame=none, xleftmargin=2.2em,
  aboveskip=6pt, belowskip=6pt,
}

\lstdefinestyle{jsonstyle}{
  basicstyle=\ttfamily\footnotesize\color{ink},
  stringstyle=\color{str},
  showstringspaces=false,
  breaklines=true,
  columns=fullflexible, keepspaces=true,
  frame=none, xleftmargin=1em,
  aboveskip=4pt, belowskip=4pt,
}

\lstdefinestyle{shellstyle}{
  basicstyle=\ttfamily\footnotesize\color{ink},
  breaklines=true, breakatwhitespace=false,
  columns=fullflexible, keepspaces=true,
  frame=none, xleftmargin=1em,
  aboveskip=4pt, belowskip=4pt,
}

%----------------------------------------------------------- boxes
% Left-ruled, page-breakable callouts built on `framed`, so the document
% compiles on a TeX Live install without the TikZ/PGF collection.
\definecolor{shadecolor}{HTML}{F4F6F7}

\newenvironment{algobox}[1]{%
  \par\addvspace{6pt}%
  \def\FrameCommand{{\color{water}\vrule width 2.5pt}\hspace{9pt}}%
  \MakeFramed{\advance\hsize-\width\FrameRestore}%
  \noindent{\small\bfseries\color{water}#1}\par\nobreak\addvspace{3pt}%
}{\par\endMakeFramed\addvspace{6pt}}

\newenvironment{notebox}[1][Note]{%
  \par\addvspace{6pt}%
  \def\FrameCommand{{\color{clay}\vrule width 2.5pt}\hspace{9pt}}%
  \MakeFramed{\advance\hsize-\width\FrameRestore}%
  \noindent{\small\bfseries\color{clay}#1}\par\nobreak\addvspace{3pt}\small%
}{\par\endMakeFramed\addvspace{6pt}}

\newenvironment{codebox}[1]{%
  \par\addvspace{6pt}%
  \def\FrameCommand{\fboxsep=8pt\colorbox{shadecolor}}%
  \MakeFramed{\FrameRestore}%
  \noindent{\ttfamily\small\color{ink}#1}\par\nobreak\addvspace{3pt}%
}{\par\endMakeFramed\addvspace{6pt}}

\newcommand{\stage}[1]{\textcolor{water}{\textbf{#1}}}

%=======================================================================
\begin{document}
%=======================================================================

%----------------------------------------------------------- title
\begin{titlepage}
\centering
\vspace*{2.2cm}

{\Huge\bfseries\color{ink} Catchment Analyzer\par}
\vspace{6mm}
{\Large\color{water} Terrain-Driven Pond Siting and\\[2mm] Catchment Estimation from Contour Maps\par}

\vspace{8mm}
{\color{rule}\rule{0.55\textwidth}{1.2pt}}
\vspace{8mm}

{\large Phase 1 --- Backend API Development\par}
\vspace{2mm}
{\normalsize A contour map goes in; a pond site and its catchment come out.\par}

\vspace{18mm}

\begin{tabular}{@{}r@{\hspace{1.2em}}l@{}}
\textbf{\color{ink}Repository}   & \url{\repourlraw} \\[4pt]
\textbf{\color{ink}API route}    & \texttt{POST \apiroute} \\[4pt]
\textbf{\color{ink}Alias route}  & \texttt{POST \apihost/api/findCatchment} \\[4pt]
\textbf{\color{ink}Live docs}    & \url{\apihost/docs} \\[4pt]
\textbf{\color{ink}Sample map}   & \texttt{contours\_1m.kml} (1355 lines, 1\,m interval) \\
\end{tabular}

\vfill
{\large Aman Pratap\par}
\vspace{1mm}
{\small Indian Institute of Technology Bhilai\par}
\vspace{4mm}
{\small\today\par}
\end{titlepage}

%----------------------------------------------------------- toc
\tableofcontents
\newpage

%=======================================================================
\section{Overview}
%=======================================================================

This report documents a backend service that accepts a contour map as
\texttt{KML} or \texttt{KMZ}, reconstructs the terrain surface it describes,
routes water across that surface, and returns --- as structured JSON --- a
recommended pond location together with the catchment that drains to it.

The central design commitment is that \emph{nothing is specific to the sample
map}. Extent, elevation range, contour interval, grid resolution and every score
threshold are derived from whatever file is uploaded. The sample map is used to
test the implementation, never to calibrate it. Section~\ref{sec:general} sets
out exactly how this is enforced.

\begin{notebox}[What the service answers]
Given only contour geometry, three questions are answered together: \emph{where}
should a pond go, \emph{how much land} drains to it, and \emph{how much water}
that land yields. A fourth is answered implicitly --- where a pond should
\emph{not} go, because the ground is a river, a floodplain, or a drainage line
too large for a farm pond to hold.
\end{notebox}

\subsection{Deliverables against the brief}

\begin{table}[H]
\centering
\small
\begin{tabular}{@{}p{0.42\textwidth}p{0.52\textwidth}@{}}
\toprule
\textbf{Requirement} & \textbf{Where it is met} \\
\midrule
\texttt{POST /analyzeContour} or \texttt{/findCatchment} & Both exist; \texttt{/findCatchment} is an alias. \S\ref{sec:api} \\
Accept contour map as uploaded file & \texttt{multipart/form-data}, KML and KMZ. \S\ref{sec:api} \\
Analyze contour / terrain information & \S\ref{sec:approach}, stages 1--4 \\
Identify a suitable pond location & \S\ref{sec:siting} \\
Estimate the catchment area & \S\ref{sec:delineate} \\
Return structured JSON & \S\ref{sec:schema} \\
Demonstrate with the sample map & \S\ref{sec:demo} \\
No hard-coded locations or results & \S\ref{sec:general} \\
Extensible to generalized maps & \S\ref{sec:extend} \\
\bottomrule
\end{tabular}
\caption{Requirement coverage.}
\end{table}

%=======================================================================
\section{System Architecture}
%=======================================================================

The service is a FastAPI application. The HTTP layer is deliberately thin: it
validates the upload and delegates immediately to a pure-Python analysis core
that has no knowledge of HTTP. That separation is what makes the core reusable
in Phase 2 from a queue worker, a batch script, or a notebook.

\subsection{Processing pipeline}

\begin{center}
\small
\begin{tabular}{@{}c@{}}
\colorbox{mist}{\parbox{0.86\textwidth}{\centering\vspace{2pt}
\stage{1. Parse} KML/KMZ $\rightarrow$ contour lines $+$ derived metadata\vspace{2pt}}}\\[3pt]
$\downarrow$\\[3pt]
\colorbox{mist}{\parbox{0.86\textwidth}{\centering\vspace{2pt}
\stage{2. Interpolate} scattered contour vertices $\rightarrow$ regular elevation grid\vspace{2pt}}}\\[3pt]
$\downarrow$\\[3pt]
\colorbox{mist}{\parbox{0.86\textwidth}{\centering\vspace{2pt}
\stage{3. Condition} priority-flood sink fill $\rightarrow$ depression depth $+$ routable surface\vspace{2pt}}}\\[3pt]
$\downarrow$\\[3pt]
\colorbox{mist}{\parbox{0.86\textwidth}{\centering\vspace{2pt}
\stage{4. Route} D8 flow direction $\rightarrow$ flow accumulation $\rightarrow$ stream network\vspace{2pt}}}\\[3pt]
$\downarrow$\\[3pt]
\colorbox{mist}{\parbox{0.86\textwidth}{\centering\vspace{2pt}
\stage{5. Classify} river / floodplain / nala / still water $\rightarrow$ exclusion mask\vspace{2pt}}}\\[3pt]
$\downarrow$\\[3pt]
\colorbox{mist}{\parbox{0.86\textwidth}{\centering\vspace{2pt}
\stage{6. Score \& rank} four weighted suitability layers $\rightarrow$ candidate sites\vspace{2pt}}}\\[3pt]
$\downarrow$\\[3pt]
\colorbox{mist}{\parbox{0.86\textwidth}{\centering\vspace{2pt}
\stage{7. Delineate} upstream walk from chosen outlet $\rightarrow$ catchment polygon\vspace{2pt}}}\\[3pt]
$\downarrow$\\[3pt]
\colorbox{mist}{\parbox{0.86\textwidth}{\centering\vspace{2pt}
\stage{8. Quantify} storage, runoff yield, time of concentration $\rightarrow$ JSON $+$ GeoJSON\vspace{2pt}}}\\
\end{tabular}
\end{center}

\subsection{Module map}

\begin{table}[H]
\centering
\small
\begin{tabular}{@{}llp{0.46\textwidth}@{}}
\toprule
\textbf{Module} & \textbf{LOC} & \textbf{Responsibility} \\
\midrule
\texttt{core/kml.py}        & 295 & Parse KML/KMZ; derive interval, bounds, levels \\
\texttt{core/grid.py}       & 183 & Interpolate contours onto a square-cell grid \\
\texttt{core/hydrology.py}  & 408 & Sink fill, D8 routing, accumulation, delineation \\
\texttt{core/siting.py}     & 518 & Suitability scoring, storage, site ranking \\
\texttt{core/watercourse.py}& 340 & River / floodplain / nala classification \\
\texttt{core/pipeline.py}   & 533 & Orchestration; assembles the response \\
\texttt{core/render.py}     & 100 & Elevation, hillshade and watercourse overlays \\
\texttt{api/routes.py}      & 209 & HTTP endpoints, validation, error mapping \\
\texttt{models.py}          & 271 & Pydantic response schema \\
\midrule
\textbf{Total} & \textbf{$\sim$3500} & \\
\bottomrule
\end{tabular}
\caption{Source layout. Each stage of the pipeline is one module with one job.}
\end{table}

%=======================================================================
\section{Catchment Estimation Approach}
\label{sec:approach}
%=======================================================================

The method is the standard raster-hydrology chain --- fill, route, accumulate,
delineate --- built on published algorithms rather than invented heuristics. What
follows is the logic in pseudocode; the repository holds the implementation.

\subsection{Stage 1 --- Contour ingestion}

A KMZ is a zip archive; the main \texttt{.kml} is extracted first. Elevation for
each placemark is taken from whichever source the file actually provides:
the coordinate $z$ component, an \texttt{<ExtendedData>} field, or a numeric
name such as \texttt{"272"}. Files in the wild use all three.

\begin{algobox}{Algorithm 1 --- Read contours from KML/KMZ}
\begin{lstlisting}[style=pseudostyle]
function read_contours(bytes, filename):
    if bytes is a zip archive then
        bytes <- extract largest .kml member   // KMZ case
    doc <- parse_xml(bytes)

    lines <- new list
    foreach placemark in doc:
        coords <- parse_coordinates(placemark)   // (lon, lat, z?) triples
        z <- first available of:
                 altitude from coordinates,
                 <ExtendedData> elevation field,
                 numeric <name> of the placemark
        if z is not null and coords has >= 2 points then
            append ContourLine(coords, z) to lines

    lines <- drop_isolated_levels(lines)   // spot heights, stray labels
    assert count(lines) >= 3 and distinct_levels(lines) >= 2

    // Everything below is DERIVED, never assumed:
    bounds   <- min/max of all vertex lon/lat
    levels   <- sorted distinct elevations
    interval <- median gap between neighbouring levels
    return ContourSet(lines, bounds, levels, interval)
\end{lstlisting}
\end{algobox}

\texttt{drop\_isolated\_levels} removes elevation values that do not belong to
the contour series --- a single placemark at 412\,m among contours spaced every
metre is a spot height, and interpolating it as terrain would raise a spurious
hill. The test is relative to the series' own median spacing, so it holds for
any interval.

\subsection{Stage 2 --- Surface reconstruction}

Contours are scattered vertices carrying a $z$ value. Turning them into a
surface is a 2-D interpolation over a Delaunay triangulation, with
nearest-neighbour fill for the region outside the convex hull, and a light
uniform filter to suppress triangulation facets.

\begin{algobox}{Algorithm 2 --- Interpolate contours onto a grid}
\begin{lstlisting}[style=pseudostyle]
function build_grid(contours, resolution = 160, margin = 0.01):
    // Pad so edge contours are not clipped
    (s, w, n, e) <- contours.bounds expanded by margin

    // Cells are kept SQUARE IN METRES -- the D8 model assumes it.
    width_m  <- haversine(mid_lat, w, mid_lat, e)
    height_m <- haversine(s, w, n, w)
    if width_m >= height_m then
        cols <- resolution;  rows <- round(resolution * height_m / width_m)
    else
        rows <- resolution;  cols <- round(resolution * width_m / height_m)
    cell_size_m <- mean(width_m / cols, height_m / rows)

    // Control points: every vertex, decimated to a bounded budget
    step <- max(1, total_vertices / MAX_CONTROL_POINTS)
    (xy, z) <- vertices[::step] with their line elevations

    targets <- regular lon/lat mesh of shape (rows, cols)   // north-up
    values  <- LinearNDInterpolator(xy, z)(targets)         // Delaunay

    outside <- cells where values is NaN                    // beyond hull
    values[outside] <- NearestNDInterpolator(xy, z)(targets[outside])

    surface <- uniform_filter(reshape(values, rows, cols), SMOOTH_WINDOW)
    return ElevationGrid(surface, bounds, cell_size_m,
                         extrapolated_fraction = mean(outside),
                         vertical_resolution_m = contours.interval)
\end{lstlisting}
\end{algobox}

\begin{notebox}[Why square cells matter]
The D8 model compares drops to eight neighbours over distances of $1$ and
$\sqrt{2}$ cell widths. If cells were rectangular those distances would be
wrong and flow would bias along one axis. Solving for the shorter side rather
than fixing both is what keeps the model valid for maps of any aspect ratio.
\end{notebox}

\subsection{Stage 3 --- Depression filling}

Routing on a raw interpolated surface strands water in every hollow, and flow
accumulation becomes meaningless. Priority-flood filling
(Wang \& Liu, 2006) raises each pit to its spill level. The fill depth added at
a cell is not a by-product to discard --- it \emph{is} the depth of the hollow
sitting there, which Stage 6 uses directly as evidence of natural storage.

\begin{algobox}{Algorithm 3 --- Priority-flood sink fill}
\begin{lstlisting}[style=pseudostyle]
function fill_sinks(z, epsilon = 0):
    filled <- copy of z
    queue  <- min-heap seeded with ALL BOUNDARY cells, keyed by elevation
    closed <- boundary cells marked

    while queue not empty:
        (level, r, c) <- pop lowest from queue
        foreach (nr, nc) in 8-neighbours of (r, c):
            if inside grid and not closed[nr, nc] then
                closed[nr, nc] <- true
                // A cell can be no lower than the lowest rim reaching it
                filled[nr, nc] <- max(filled[nr, nc], level + epsilon)
                push (filled[nr, nc], nr, nc) onto queue
    return filled

function depression_depth(z):
    return max(0, fill_sinks(z, 0) - z)     // head of water standing per cell

function conditioned_surface(z):
    return fill_sinks(z, epsilon = 1e-4)    // tilt flats so flow is defined
\end{lstlisting}
\end{algobox}

The $\epsilon$ tilt matters. A filled pit is a dead-flat lake in which no cell
has a downhill neighbour, so flow direction is undefined across it. Adding
$10^{-4}$\,m per step of the flood restores a drainage path without changing any
elevation meaningfully (Barnes et al., 2014).

\subsection{Stage 4 --- D8 flow routing}

\begin{algobox}{Algorithm 4 --- Flow direction and accumulation}
\begin{lstlisting}[style=pseudostyle]
function flow_direction(z, cell_size_m):
    // D8: each cell drains to its steepest downhill neighbour of eight.
    padded <- z padded by one ring of +infinity   // off-grid can never win
    best   <- zeros;  direction <- filled with -1 // -1 marks a pit

    foreach k, (dr, dc) in the 8 offsets:
        neighbour <- padded shifted by (dr, dc)
        slope     <- (z - neighbour) / (cell_size_m * dist[k])
        steeper   <- slope > best
        best      <- where(steeper, slope, best)
        direction <- where(steeper, k, direction)
    return direction

function flow_accumulation(z, direction):
    // Water only moves downhill, so processing cells from HIGHEST to LOWEST
    // guarantees a cell's total is final before it is passed on.
    accumulation <- ones(size of z)             // each cell counts itself
    foreach cell in cells sorted by z descending:
        d <- downstream neighbour of cell via direction
        if d exists then
            accumulation[d] <- accumulation[d] + accumulation[cell]
    return accumulation
\end{lstlisting}
\end{algobox}

Sorting by elevation replaces the usual recursive traversal with a single
linear pass, which is what keeps a $131 \times 160$ grid inside half a second.

\subsection{Stage 5 --- What the ground already is}

Before any site is proposed, the map is split into ground a farm pond may
occupy and ground it may not. The split is by \emph{upstream area in hectares},
not by any property of this particular map, so it transfers unchanged to other
terrain.

\begin{table}[H]
\centering
\small
\begin{tabular}{@{}lll@{}}
\toprule
\textbf{Class} & \textbf{Test} & \textbf{Consequence} \\
\midrule
Farm pond ground & upstream $\le 25$\,ha            & pond proposed \\
Nala             & $25 < $ upstream $\le 100$\,ha   & bund / percolation tank \\
River            & upstream $> 100$\,ha             & no structure \\
Floodplain       & within 40\,m buffer, HAND $< 3$\,m & no structure \\
Still water      & flat pooled area $\ge 2$\,ha     & no structure \\
\bottomrule
\end{tabular}
\caption{Classification thresholds. All are physical, in hectares or metres.}
\end{table}

\begin{notebox}[The reasoning behind 25 ha]
ICAR dryland manuals and the MGNREGA works schedule put a farm pond's catchment
at roughly 1--10\,ha. Past about 25\,ha an ordinary storm delivers more than the
pond can retain, and the engineering problem becomes a spillway rather than an
embankment --- a different structure, correctly reported as such rather than
silently mis-recommended.
\end{notebox}

\subsection{Stage 6 --- Pond siting}
\label{sec:siting}

Every eligible cell is scored on four things a contour map can genuinely answer.
Each layer is normalised against the range present \emph{in the uploaded map},
so the score is relative to that terrain and never to an absolute constant.

\begin{align}
n(x) &= \frac{x - \min x}{\max x - \min x} \\[4pt]
L_{\text{depression}} &= n(\text{fill depth}) \\
L_{\text{catchment}}  &= n\!\left(\log(1 + A)\right) \\
L_{\text{slope}}      &= n\!\left(\frac{1}{1 + \theta / \theta_{1/2}}\right),
                         \quad \theta_{1/2} = 8^{\circ} \\
L_{\text{elevation}}  &= n(-z) \\[6pt]
S &= 100 \cdot \frac{\sum_i w_i L_i}{\sum_i w_i}
\end{align}

with weights $w = \{$depression $0.30$, catchment $0.30$, slope $0.25$,
elevation $0.15\}$. The $\log$ on accumulation matters: raw upstream counts span
orders of magnitude, and without it a single trunk channel would saturate the
layer and flatten every genuine difference elsewhere.

\begin{algobox}{Algorithm 5 --- Score, filter and rank candidate sites}
\begin{lstlisting}[style=pseudostyle]
function select_sites(grid, accumulation, depths, direction, max_sites):
    layers <- { depression: n(depths),
                catchment:  n(log1p(accumulation)),
                slope:      n(1 / (1 + slope_deg / 8)),
                elevation:  n(-grid.z) }
    score  <- 100 * weighted_sum(layers, WEIGHTS) / sum(WEIGHTS)

    // Rule out ground that cannot take a pond
    blocked <- river or floodplain or still_water or edge_ring
    // A busy channel with no measurable hollow is interpolation, not storage
    holds   <- max(0.30, grid.vertical_resolution_m)
    blocked <- blocked or (accumulation >= percentile(accumulation, 95)
                           and depths < holds)

    ranked <- cells not blocked, sorted by score descending
    sites  <- new list
    foreach cell in ranked:
        if count(sites) = max_sites then break
        if not spaced_apart(sites, cell, separation) then continue

        storage <- natural_storage(grid, depths, cell)
        inflow  <- upstream_cells * cell_area * 0.4 * 50mm / 1000  // one storm
        if inflow > storage.volume then
            reclassify as channel structure (bund + waste weir)
        else
            append PondSite(cell, score, storage, reasons) to sites
    return sites
\end{lstlisting}
\end{algobox}

Storage is not assumed from a design depth. The sink fill already gives the
spill elevation, so the pond is grown outwards from the site up to that real rim:

\begin{algobox}{Algorithm 6 --- Natural storage by flood fill to the spill rim}
\begin{lstlisting}[style=pseudostyle]
function natural_storage(grid, depths, cell):
    if depths[cell] <= 0 then return zero storage    // nothing to hold; must be dug
    spill <- grid.z[cell] + depths[cell]             // the real rim it overtops

    flooded <- {cell};  queue <- [cell];  head_sum <- 0
    while queue not empty:
        (r, c) <- pop front
        head_sum <- head_sum + (spill - grid.z[r, c])
        foreach 4-neighbour (nr, nc) not already flooded:
            if grid.z[nr, nc] < spill then
                add (nr, nc) to flooded;  push onto queue

    return { depth_m:        depths[cell],
             spill_m:        spill,
             surface_area:   count(flooded) * cell_area,
             volume_m3:      head_sum * cell_area }
\end{lstlisting}
\end{algobox}

\subsection{Stage 7 --- Catchment delineation}
\label{sec:delineate}

With the site chosen, its catchment is exactly the set of cells whose flow
pointers lead to it --- found by walking the D8 network \emph{upstream} from the
outlet.

\begin{algobox}{Algorithm 7 --- Delineate the catchment of an outlet}
\begin{lstlisting}[style=pseudostyle]
function delineate(grid, direction, outlet):
    contributors <- invert(direction)      // downstream pointer -> upstream list
    cells <- {outlet};  queue <- [outlet]
    while queue not empty:
        cell <- pop front
        foreach up in contributors[cell]:
            if up not in cells then
                add up to cells;  push up onto queue

    area_m2   <- count(cells) * grid.cell_area_m2
    polygon   <- union of cell squares, exterior ring simplified
    flow_path <- longest upstream chain from outlet, in metres
    gradient  <- (max_z - min_z) / flow_path
    return Catchment(cells, area_m2, polygon, flow_path, gradient, relief, slope)
\end{lstlisting}
\end{algobox}

\subsection{Stage 8 --- Quantifying yield}

Two standard formulae close the analysis. The Rational method gives the water a
catchment yields; Kirpich gives how quickly it arrives, which is what sizes a
spillway.

\begin{align}
V &= P \cdot A \cdot C
   &&\text{runoff volume (m\textsuperscript{3})} \\
t_c &= 0.0195\, L^{0.77} S^{-0.385}
   &&\text{time of concentration (min)}
\end{align}

where $P$ is rainfall depth (mm), $A$ catchment area (m\textsuperscript{2}),
$C$ the runoff coefficient (default $0.4$), $L$ the longest flow path (m) and
$S$ the average gradient. When no rainfall is supplied the response reports the
yield per millimetre, $A \cdot C / 1000$, so the caller may apply local rainfall
themselves.

%=======================================================================
\section{Generalization: no hard-coded results}
\label{sec:general}
%=======================================================================

This is the requirement most easily failed by accident, so it is worth being
explicit about what is derived and what is constant.

\begin{table}[H]
\centering
\small
\begin{tabular}{@{}p{0.34\textwidth}p{0.28\textwidth}p{0.30\textwidth}@{}}
\toprule
\textbf{Quantity} & \textbf{Source} & \textbf{Sample value} \\
\midrule
Geographic extent      & vertex min/max        & 21.240--21.264\,N \\
Elevation range        & contour $z$ values    & 267--298\,m \\
Contour interval       & median level gap      & 1.0\,m \\
Grid rows $\times$ cols& extent aspect ratio   & $131 \times 160$ \\
Cell size              & haversine / cols      & 20.6\,m \\
Score normalisation    & map's own min/max     & per-upload \\
Pond location          & argmax of score       & computed \\
Catchment area         & upstream cell count   & computed \\
\bottomrule
\end{tabular}
\caption{Every map-specific quantity is derived at request time.}
\end{table}

The only constants in the codebase are physical or engineering ones that hold
anywhere: the 25\,ha and 100\,ha structure thresholds, the $8^{\circ}$
slope half-score, the 50\,mm storm test, Kirpich's coefficients, and the D8
neighbour geometry. None encodes a location.

\begin{notebox}[A concrete check]
Uploading a different contour map produces a different grid shape, a different
cell size, different normalisation ranges and a different site --- with no code
change. The companion endpoint \texttt{/analyzeArea} exercises the same core on
terrain downloaded for arbitrary world coordinates, which is the strongest
available evidence that the analysis is not tuned to the sample.
\end{notebox}

%=======================================================================
\section{API Documentation}
\label{sec:api}
%=======================================================================

\subsection{Endpoints}

\begin{table}[H]
\centering
\small
\begin{tabular}{@{}llp{0.40\textwidth}@{}}
\toprule
\textbf{Method} & \textbf{Path} & \textbf{Purpose} \\
\midrule
\texttt{POST} & \texttt{/api/analyzeContour} & Analyse an uploaded contour map \\
\texttt{POST} & \texttt{/api/findCatchment}  & Alias of the above \\
\texttt{POST} & \texttt{/api/analyzeArea}    & Analyse by coordinates (no file) \\
\texttt{GET}  & \texttt{/api/elevation}      & Point elevation lookup \\
\texttt{GET}  & \texttt{/api/health}         & Liveness check \\
\texttt{GET}  & \texttt{/docs}               & Interactive OpenAPI documentation \\
\bottomrule
\end{tabular}
\caption{Both prefixed (\texttt{/api/...}) and bare (\texttt{/...}) paths are served.}
\end{table}

\begin{notebox}[The analysis routes are POST-only]
\texttt{/api/analyzeContour} carries a file in the request body and therefore
cannot be opened from a browser address bar --- a \texttt{GET} correctly returns
\texttt{405 Method Not Allowed}. To exercise it interactively, open
\url{\apihost/docs}, which provides an upload form bound to the live endpoint.
\end{notebox}

\subsection{Request}

\texttt{Content-Type: multipart/form-data}

\begin{table}[H]
\centering
\small
\begin{tabular}{@{}lllp{0.34\textwidth}@{}}
\toprule
\textbf{Field} & \textbf{Type} & \textbf{Default} & \textbf{Meaning} \\
\midrule
\texttt{file}              & file  & required & Contour map, \texttt{.kml} or \texttt{.kmz}, $\le 64$\,MB \\
\texttt{resolution}        & int   & 160      & Cells along the longer side (48--260) \\
\texttt{max\_sites}        & int   & 5        & Candidate sites to return \\
\texttt{runoff\_coefficient}& float& 0.4      & $C$ in the Rational method \\
\texttt{rainfall\_mm}      & float & null     & If given, absolute runoff volume is computed \\
\bottomrule
\end{tabular}
\caption{Request parameters. Only \texttt{file} is mandatory.}
\end{table}

\subsection{Response schema}
\label{sec:schema}

\begin{codebox}{200 OK --- abridged}
\begin{lstlisting}[style=jsonstyle]
{
  "success": true,
  "analysis_id": "ca_aff346d081",
  "generated_at": "2026-09-01T17:28:13+00:00",
  "source":   { "kind", "format", "contour_lines", "vertices",
                "elevation_levels", "contour_interval_m", "bounds" },
  "terrain":  { "min_elevation_m", "max_elevation_m", "relief_m",
                "mean_slope_deg", "mapped_area_km2", "grid" },
  "pond_site":{ "rank", "latitude", "longitude", "elevation_m",
                "slope_deg", "depression_depth_m", "upstream_hectares",
                "height_above_drainage_m", "structure", "score",
                "rating", "score_breakdown", "storage", "reasons" },
  "catchment":{ "outlet", "area_m2", "area_km2", "area_hectares",
                "cell_count", "perimeter_m", "relief_m",
                "longest_flow_path_m", "average_gradient",
                "time_of_concentration_min", "runoff" },
  "alternative_sites":  [ ... ],
  "channel_structures": [ ... ],
  "watercourses":       { ... },
  "overlays":           { "elevation", "hillshade", "watercourse" },
  "geojson":            { "type": "FeatureCollection", "features": [ ... ] },
  "warnings":           [ ... ]
}
\end{lstlisting}
\end{codebox}

\subsection{Error responses}

\begin{table}[H]
\centering
\small
\begin{tabular}{@{}llp{0.44\textwidth}@{}}
\toprule
\textbf{Status} & \textbf{Meaning} & \textbf{Typical cause} \\
\midrule
400 & Bad request      & Not a KML/KMZ, or unreadable archive \\
405 & Method not allowed & \texttt{GET} on a POST-only analysis route \\
413 & Payload too large& Upload exceeds 64\,MB \\
422 & Unprocessable    & Fewer than 3 contour lines or 2 distinct levels \\
500 & Analysis failed  & Interpolation produced an incomplete surface \\
\bottomrule
\end{tabular}
\caption{Errors share one shape: \texttt{\{success, error, detail\}}.}
\end{table}

%=======================================================================
\section{Demonstration}
\label{sec:demo}
%=======================================================================

Run against the provided \texttt{contours\_1m.kml}:

\begin{codebox}{Request}
\begin{lstlisting}[style=shellstyle]
curl -X POST \apiroute \
     -F "file=@contours_1m.kml" \
     -F "resolution=160"
\end{lstlisting}
\end{codebox}

\noindent Observed: \texttt{HTTP 200}, 43{,}588 bytes, in approximately
0.5\,s.

\subsection{What was read from the file}

\begin{table}[H]
\centering
\small
\begin{tabular}{@{}lr@{}}
\toprule
\textbf{Property} & \textbf{Derived value} \\
\midrule
Contour lines            & 1{,}355 \\
Vertices                 & 159{,}113 \\
Distinct elevation levels& 32 \\
Contour interval         & 1.0\,m \\
Elevation range          & 267.0 -- 298.0\,m \\
Bounds (S, W, N, E)      & 21.2398, 81.2814, 21.2636, 81.3126 \\
\bottomrule
\end{tabular}
\caption{Stage 1 output --- entirely derived from the upload.}
\end{table}

\subsection{Reconstructed terrain}

\begin{table}[H]
\centering
\small
\begin{tabular}{@{}lr@{}}
\toprule
Grid                     & $131 \times 160$ cells \\
Cell size / area         & 20.605\,m / 424.57\,m\textsuperscript{2} \\
Mapped area              & 8.8989\,km\textsuperscript{2} \\
Elevation (min/mean/max) & 267.55 / 283.83 / 296.67\,m \\
Relief                   & 29.13\,m \\
Mean / max slope         & $2.37^{\circ}$ / $13.36^{\circ}$ \\
Extrapolated fraction    & 0.0755 \\
\bottomrule
\end{tabular}
\caption{Stage 2--3 output.}
\end{table}

\subsection{Recommended pond site and its catchment}

\begin{table}[H]
\centering
\small
\begin{tabular}{@{}lr@{}}
\toprule
\multicolumn{2}{@{}l}{\textbf{\color{water}Pond site (rank 1)}} \\
\midrule
Location                 & 21.249651\,N, 81.286503\,E \\
Structure                & Farm pond \\
Score / rating           & 54.3 --- Moderate \\
Ground elevation         & 277.38\,m \\
Slope                    & $1.11^{\circ}$ \\
Depression depth         & 2.74\,m \\
Height above drainage    & 5.66\,m \\
Spill elevation          & 280.12\,m \\
Water surface area       & 21{,}228.3\,m\textsuperscript{2} \\
\textbf{Storage volume}  & \textbf{26{,}811.3\,m\textsuperscript{3}} \\
\midrule
\multicolumn{2}{@{}l}{\textbf{\color{water}Catchment}} \\
\midrule
Outlet                   & site coordinates \\
\textbf{Area}            & \textbf{9.47\,ha} (94{,}678.2\,m\textsuperscript{2}) \\
Cell count               & 223 \\
Perimeter                & 1{,}867.6\,m \\
Elevation range          & 275.66 -- 285.27\,m (relief 9.61\,m) \\
Mean / max slope         & $2.84^{\circ}$ / $7.48^{\circ}$ \\
Longest flow path        & 601.4\,m \\
Average gradient         & 0.01599 \\
Time of concentration    & 13.2\,min \\
Runoff yield             & 37.87\,m\textsuperscript{3} per mm of rainfall \\
Share of mapped area     & 1.06\,\% \\
\bottomrule
\end{tabular}
\caption{Primary result for the sample contour map.}
\end{table}

\subsection{Arithmetic check}

Both reported quantities reproduce from the formulae in \S\ref{sec:approach},
which is a useful guard against a plausible-looking but wrong number:

\begin{align}
\text{yield} &= \frac{A \cdot C}{1000}
              = \frac{94678.2 \times 0.4}{1000}
              = 37.87\ \text{m\textsuperscript{3}/mm} \quad\checkmark \\
t_c &= 0.0195 \times 601.4^{0.77} \times 0.01599^{-0.385}
     = 13.2\ \text{min} \quad\checkmark
\end{align}

\subsection{Ground the analysis refused}

The map is not uniformly available for a farm pond, and the service says so
rather than proposing a site in a river:

\begin{table}[H]
\centering
\small
\begin{tabular}{@{}lr@{}}
\toprule
River                & 51.46\,ha \\
Floodplain           & 229.18\,ha \\
Nala                 & 9.38\,ha \\
Truncated channels   & 1.70\,ha \\
Still water          & 0.00\,ha \\
\midrule
\textbf{Excluded fraction} & \textbf{32.59\,\%} \\
\bottomrule
\end{tabular}
\caption{Stage 5 output. Roughly a third of the map carries no pond proposal.}
\end{table}

Where a pond is inappropriate the service proposes the structure that
\emph{is} appropriate rather than returning nothing. The best such alternative
on this map is a nala bund at 21.247414\,N, 81.289911\,E, serving 51.46\,ha,
scoring 65.5 (\emph{Good}) --- a higher score than the pond, correctly reported
as a different structure.

\subsection{Map output}

Alongside the JSON, the response carries GeoJSON (12 features: catchment
polygon, site markers, stream network, longest flow path) and three PNG
overlays --- elevation, hillshade and watercourse classification --- georeferenced
to the analysis bounds.

\IfFileExists{figures/catchment-map.png}{%
  \begin{figure}[H]
  \centering
  \includegraphics[width=\textwidth]{figures/catchment-map.png}
  \caption{Analysis of the sample map. Marker \textbf{1} is the recommended
  pond site with its catchment outlined; \textbf{2}--\textbf{5} are ranked
  alternatives; \textbf{B1}--\textbf{B3} are nala-bund proposals on drainage
  lines too large for a farm pond. The blue trunk is the classified river.}
  \label{fig:map}
  \end{figure}
}{%
  \begin{notebox}[Figure placeholder]
  Save the map screenshot as \texttt{report/figures/catchment-map.png} and
  recompile; the figure and its caption will appear here automatically.
  \end{notebox}
}

%=======================================================================
\section{Extensibility to Phase 2}
\label{sec:extend}
%=======================================================================

The architecture anticipates generalization in four specific ways.

\begin{enumerate}[label=\textbf{\arabic*.}]
\item \textbf{Terrain source is pluggable.} \texttt{build\_grid} and
\texttt{grid\_shape} are shared by the contour path and the DEM-download path.
A grid built from a downloaded raster has exactly the geometry one built from
contours would, so every downstream stage is source-agnostic. Adding GeoTIFF or
Shapefile input means writing one reader that returns a \texttt{ContourSet} ---
nothing in hydrology, siting or the API changes.

\item \textbf{Thresholds are physical, not positional.} Structure
classification is in hectares and metres. Moving to a different district,
state or country requires no retuning.

\item \textbf{Resolution is a request parameter.} The same map can be analysed
coarsely for a fast preview or finely for a design study (48--260 cells per
side) without redeployment.

\item \textbf{The core is HTTP-free.} \texttt{core/} imports nothing from
FastAPI. The identical pipeline can be driven from a batch job, a scheduled
task or a notebook in Phase 2.
\end{enumerate}

\subsection{Known limitations}

Stated plainly, since they bound what the current results mean:

\begin{itemize}
\item Interpolation between contours cannot recover features finer than the
contour interval; on this map, hollows shallower than 1\,m are not
distinguishable from smoothing artefacts, which is why the channel test uses
the interval as its floor.
\item About 7.6\,\% of the sample grid lies outside the contour convex hull and
is nearest-neighbour extrapolated; the outer ring is excluded from siting for
this reason.
\item Catchments are clipped by the map extent. Flow entering from beyond the
mapped area is not counted, and the response emits a warning when a candidate's
catchment touches the boundary.
\item Runoff uses the Rational method, which is appropriate for small
catchments ($<$ 25\,ha here) but not for basin-scale routing.
\end{itemize}

%=======================================================================
\section{Reproducing the result}
%=======================================================================

\begin{codebox}{Local}
\begin{lstlisting}[style=shellstyle]
git clone \repourlraw
cd catchment_analyzer
python3 -m venv .venv && .venv/bin/pip install -r backend/requirements.txt
.venv/bin/uvicorn backend.main:app --host 0.0.0.0 --port 5249

curl -X POST http://localhost:5249/api/analyzeContour -F "file=@contours_1m.kml"
\end{lstlisting}
\end{codebox}

\begin{codebox}{Docker}
\begin{lstlisting}[style=shellstyle]
docker compose up --build      # API and UI published on port 5249
\end{lstlisting}
\end{codebox}

%=======================================================================
\section*{References}
\addcontentsline{toc}{section}{References}
%=======================================================================

\begin{enumerate}[label={[\arabic*]}]
\item O'Callaghan, J.F. \& Mark, D.M. (1984). The extraction of drainage
networks from digital elevation data. \emph{Computer Vision, Graphics and
Image Processing}, 28(3), 323--344.
\item Wang, L. \& Liu, H. (2006). An efficient method for identifying and
filling surface depressions in digital elevation models.
\emph{International Journal of Geographical Information Science}, 20(2), 193--213.
\item Barnes, R., Lehman, C. \& Mulla, D. (2014). Priority-flood: an optimal
depression-filling and watershed-labeling algorithm for digital elevation
models. \emph{Computers \& Geosciences}, 62, 117--127.
\item Kirpich, Z.P. (1940). Time of concentration of small agricultural
watersheds. \emph{Civil Engineering}, 10(6), 362.
\item Indian Council of Agricultural Research. \emph{Dryland Agriculture
Manual --- Farm Pond Design}.
\end{enumerate}

\end{document}
